\input{./._relpath-to-latexroot.ltx}
\documentclass[\latexroot/\projectname]{subfiles}
\input{\latexroot/@resources/subfile-setup.ltx}

\begin{document}

\section{Model}\whenintegrated{\label{model}}
\whenintegrated{\label{sec:model}} % integrated doc: SST for labels

The model is a small open economy with nominal price rigidities as in \cite{GaliMonacelli2005}, augmented with capital and domestically incomplete markets for households as in \cite{Huggett1993,Aiyagari1994}.

\subsection{Households}\whenintegrated{\label{sec:households}}

The economy consists of a continuum of agents normalized to measure~1 who are ex-ante homogeneous, and have CRRA preferences over consumption of home and foreign goods and leisure:
\begin{equation}
  U = E_0 \sum_{t=0}^{\infty} \beta^t \left[ u(c_t) - g(l_t) \right]
  \whenintegrated{\label{eq:utility}}
\end{equation}
where $\beta \in (0,1)$ is the households' subjective discount factor, $g(l_t) = \varphi \frac{l_t^{1+1/\eta}}{1+1/\eta}$ represents the disutility of supplying labor $l_t$, and $c_t$ is a CES aggregate of home and foreign goods with elasticity of substitution $\theta$.

Households' labor productivity $\{s_t\}_{t=0}^{\infty}$ is stochastic and is characterized by an $N$-state Markov chain that can take values $s_t \in S = \{s_1, \dots, s_N\}$ with transition probabilities given by $\gamma(s_{t+1} \mid s_t)$.

At time $t$ an agent faces the following nominal budget constraint:
\begin{equation}
  P_t c_t + a_{t+1} = (1 + i_t^a) a_t + W_t l_t s_t + d_t
  \whenintegrated{\label{eq:budget}}
\end{equation}
where $P_t$ is the consumer price index, $d_t$ are dividends distributed by firms, and $i_t^a$ is the nominal return on the asset portfolio. In addition, households are subject to a borrowing constraint $a_{t+1} \geq \bar{a}$.

The household's problem can be written recursively as:
\begin{equation}
  V(a,s;\Omega)
  =
  \max_{c \geq 0,\, a' \geq \bar{a}}
  \left\{
  u(c) - g(l)
  +
  \beta
  \sum_{s' \in S}
  \gamma(s' \mid s)\,
  V(a', s'; \Omega')
  \right\}
  \whenintegrated{\label{eq:bellman}}
\end{equation}
subject to $Pc + a' = (1 + i^a)a + Wls + \pi$ and $\Omega' = \Upsilon(\Omega)$, where $\Omega$ is the distribution of agents' asset holdings and labor productivity across the population.

\subsection{Households' portfolio}\whenintegrated{\label{sec:portfolio}}

There are two broad classes of assets, bonds and capital. Nominal assets carried by the household into the following period are:
\begin{equation}
  a_{t+1} = b_{H,t+1} + e_t b_{F,t+1} + P_{H,t}(q_t k_{t+1} + \tilde{b}_{H,t+1})
  \whenintegrated{\label{eq:portfolio}}
\end{equation}
where $b_{H,t+1}$ and $b_{F,t+1}$ denote nominal bonds in domestic and foreign currency, $e_t$ is the nominal exchange rate, $q_t$ is the real price of capital, and $\tilde{b}_{H,t+1}$ are real bonds in units of home good.

\subsection{Production}\whenintegrated{\label{sec:production}}

A representative final-goods firm operates a CES production technology that aggregates domestic intermediate goods. Intermediate-goods firms operate a Cobb-Douglas technology $Y_{jt} = K_{jt}^\alpha L_{jt}^{1-\alpha}$ and set prices subject to Rotemberg adjustment costs, yielding a standard New Keynesian Phillips Curve.

\subsection{Monetary policy}\whenintegrated{\label{sec:monetary}}

We consider two regimes: (i)~a fixed exchange rate regime where $e_t = \bar{e}$ for all $t$, and (ii)~a flexible exchange rate regime where the central bank adjusts the exchange rate to achieve the flexible price allocation.

% Smart bibliography: Only include bibliography if standalone AND has citations
\smartbib

\end{document}
